\chapter{June Huh --- Fields Medal 2022}\label{ch2022junehuh:june-huh-fields-medal-2022-a-comprehensive-report-on-his-work-and-impact}

\begin{center}\rule{0.5\linewidth}{0.5pt}\end{center}

\textbf{Prepared:} July 2026 \textbf{Subject:} June E Huh (b. June 9,
1983, Stanford, California, USA) \textbf{Position:} Professor, Princeton
University \textbf{Doctoral Advisor:} Mircea Mustata \textbf{Fields
Medal Citation:} \emph{``For bringing the ideas of Hodge theory to
combinatorics, the proof of the Dowling--Wilson conjecture for geometric
lattices, the proof of the Heron--Rota--Welsh conjecture for matroids,
the development of the theory of Lorentzian polynomials, and the proof
of the strong Mason conjecture.''}

\begin{center}\rule{0.5\linewidth}{0.5pt}\end{center}

\section{Table of Contents}\label{ch2022junehuh:table-of-contents}

\begin{enumerate}
\def\labelenumi{\arabic{enumi}.}
\tightlist
\item
 Introduction
\item
 Biographical Background: An Unconventional Path to Mathematics
\item
 The Chromatic Polynomial and the Read Conjecture
\item
 The Heron--Rota--Welsh Conjecture
\item
 Hodge Theory for Matroids
\item
 The Dowling--Wilson Top-Heavy Conjecture
\item
 Lorentzian Polynomials
\item
 The Strong Mason Conjecture
\item
 Lagrangian Geometry of Matroids
\item
 Stellahedral Geometry of Matroids
\item
 The Semi-Small Decomposition and Intersection Cohomology of Matroids
\item
 Combinatorial Applications of the Hodge--Riemann Relations
\item
 Broader Impact on Mathematics
\item
 Unusual Career Trajectory: Lessons for the Profession
\item
 Awards and Recognition
\item
 Selected Major Publications
\item
 Ongoing Work and Future Directions
\item
 Conclusion
\item
 References and Further Reading
\end{enumerate}

\begin{center}\rule{0.5\linewidth}{0.5pt}\end{center}

\section{1. Introduction}\label{ch2022junehuh:introduction}

June Huh was awarded the Fields Medal in 2022 for a body of work that
represents a notable cross-fertilization in modern
mathematics: the importation of deep geometric ideas from Hodge theory
into the discrete world of combinatorics. His proofs of the
Heron--Rota--Welsh conjecture, the Dowling--Wilson conjecture, and the
strong Mason conjecture resolved problems that had stood open for
decades, and his development of the theory of Lorentzian polynomials has
provided the field with new tools whose reach extends far
beyond their original context.

A notable aspect of his personal
trajectory. Unlike many Fields Medalists who showed exceptional
mathematical promise from an early age, Huh struggled with mathematics
in school, dropped out of high school to write poetry, and only
discovered his passion for mathematics in his mid-twenties after a
chance encounter with a noted mathematician. His story has become an
inspiration to students and researchers worldwide, demonstrating that
mathematical creativity can flourish at any stage of life.

This report provides a detailed account of Huh's important
contributions, the mathematical context in which they arose, the
techniques he and his collaborators developed, and their lasting impact
on mathematics.

\begin{center}\rule{0.5\linewidth}{0.5pt}\end{center}

\section{2. Biographical Background: An Unconventional Path to
Mathematics}\label{ch2022junehuh:biographical-background-an-unconventional-path-to-mathematics}

\subsection{2.1 Early Life}\label{ch2022junehuh:early-life}

June E Huh was born on June 9, 1983, in Stanford, California, while his
Korean parents were completing graduate studies at Stanford University.
His father later became a professor of statistics at Korea University,
and his mother a professor of Russian language at Seoul National
University. When Huh was about two years old, the family returned to
South Korea, where he was raised and educated.

Huh's early experience with mathematics was discouraging. Poor scores on
elementary school tests convinced him that he lacked the innate aptitude
for the subject. He later described himself as being ``pretty good at
most subjects except math.'' In high school, he became bored and
exhausted by the relentless routine of studying and chose to drop out to
focus on writing poetry. For a period, he considered becoming a science
journalist or a poet.

\subsection{2.2 Undergraduate Studies}\label{ch2022junehuh:undergraduate-studies}

In 2002, Huh matriculated at Seoul National University (SNU). He
initially felt unsettled and suffered from depression. He pinned his
career aspirations on science journalism and decided to major in physics
and astronomy. However, his attendance was poor, and he had to repeat
several courses he initially failed.

The turning point came in his final year of undergraduate study. He
noticed that the notable Japanese mathematician Heisuke Hironaka --- a
Fields Medalist from 1970 --- was visiting SNU and giving a lecture
course on algebraic geometry. Huh knew nothing about the mathematics,
but he had read Hironaka's autobiography and found him fascinating. He
enrolled in the course.

\subsection{2.3 The Encounter with
Hironaka}\label{ch2022junehuh:the-encounter-with-hironaka}

Hironaka's course was unconventional. Rather than presenting polished,
textbook material, he used the lectures to report on his current
research, often describing insights he had literally had the day before.
For Huh, this was revelatory: ``This was the first time I actually saw
someone doing maths\ldots{} It was my first exposure to mathematics as a
human activity.''

Huh was captivated by the creative, exploratory nature of mathematical
research. He began meeting regularly with Hironaka, who became his
mentor. Hironaka encouraged Huh to pursue graduate studies in
mathematics and helped arrange his admission to graduate programs.

\subsection{2.4 Graduate Studies and
PhD}\label{ch2022junehuh:graduate-studies-and-phd}

Huh completed a master's degree in mathematics at Seoul National
University, then spent time at the University of Illinois at
Urbana-Champaign before enrolling in the PhD program at the University
of Michigan. His doctoral advisor was Mircea Mustata, a leading
algebraic geometer.

Huh's PhD thesis, completed in 2014 and titled \emph{``Rota's conjecture
and positivity of algebraic cycles in permutohedral varieties,''}
already contained the seeds of his important contributions. In it,
he proved the Read conjecture on the unimodality of chromatic polynomial
coefficients --- the first major step toward the resolution of the
Heron--Rota--Welsh conjecture.

\subsection{2.5 Academic Career}\label{ch2022junehuh:academic-career}

After his PhD, Huh held positions at the Institute for Advanced Study
(as a Clay Fellow), and then as a professor at Stanford University. In
2021, he moved to Princeton University, where he is currently a
professor. He spends his summers at the Korea Institute for Advanced
Study.

\begin{center}\rule{0.5\linewidth}{0.5pt}\end{center}

\section{3. The Chromatic Polynomial and the Read
Conjecture}\label{ch2022junehuh:the-chromatic-polynomial-and-the-read-conjecture}

\subsection{3.1 Background}\label{ch2022junehuh:background}

The chromatic polynomial \(\chi_G(\lambda)\) of a graph \(G\) is a
central object in algebraic combinatorics. It counts the number of
proper colorings of the vertices of \(G\) using \(\lambda\) colors,
where adjacent vertices must receive different colors. The polynomial
was introduced by George Birkhoff in 1912 as a tool for studying the
four-color problem.

Written in expanded form:

\[
\chi_G(\lambda) = \sum_{i=0}^n (-1)^i a_i \lambda^{n-i}
\]

the coefficients \(a_i\) encode deep information about the graph's
structure. A long-standing conjecture, made by Ronald Read in 1968,
asserted that the absolute values of the coefficients \(|a_i|\) form a
\textbf{unimodal} sequence:
\(|a_0| \le |a_1| \le \cdots \le |a_{k}| \ge \cdots \ge |a_n|\) for some
\(k\). Even more strongly, the sequence was conjectured to be
\textbf{log-concave}: \(a_i^2 \ge a_{i-1} a_{i+1}\) for all \(i\).

\subsection{3.2 Huh's Proof}\label{ch2022junehuh:huhs-proof}

In his first major result (published in the \emph{Journal of the
American Mathematical Society} in 2012), Huh proved Read's conjecture.
The proof was striking because it used techniques from \textbf{algebraic
geometry} --- specifically, the theory of \textbf{Milnor numbers} of
isolated hypersurface singularities --- to solve a purely combinatorial
problem.

The key idea was to associate to a graph \(G\) a projective hypersurface
with isolated singularities whose Milnor numbers are precisely the
coefficients of the chromatic polynomial (up to sign). The Milnor
numbers of such a hypersurface satisfy constraints (the Monodromy
conjecture, proved by Varchenko and others) that imply the desired
unimodality and log-concavity.

More precisely, Huh showed that the coefficients \(a_i\) can be
interpreted as the dimensions of certain cohomology groups associated to
a \textbf{hyperplane arrangement} derived from the graph. The hard
Lefschetz theorem for the cohomology of these spaces then forces the
coefficients to be log-concave.

\subsection{3.3 Impact}\label{ch2022junehuh:impact}

This result was a significant result. It was the first time that deep theorems
from algebraic geometry had been used to settle a purely combinatorial
conjecture about graphs. It demonstrated that the chromatic polynomial,
which had been studied extensively through combinatorial methods, had a
hidden geometric structure that could be exploited in unexpected ways.

The proof also revealed a connection between graph theory and the theory
of hyperplane arrangements that would prove crucial in Huh's later work.

\begin{center}\rule{0.5\linewidth}{0.5pt}\end{center}

\section{4. The Heron--Rota--Welsh
Conjecture}\label{ch2022junehuh:the-heronrotawelsh-conjecture}

\subsection{4.1 Background}\label{ch2022junehuh:background-1}

Matroids are a fundamental abstraction in combinatorics, generalizing
the notion of linear independence from linear algebra and the notion of
cycles from graph theory. A matroid \(M\) on a finite set \(E\) is
defined by a rank function that satisfies certain axioms; the
\textbf{characteristic polynomial} of \(M\) is defined analogously to
the chromatic polynomial of a graph.

The Heron--Rota--Welsh conjecture (also called the Rota conjecture)
states that for any matroid \(M\), the absolute values of the
coefficients of its characteristic polynomial form a
\textbf{log-concave} sequence. This was proposed by Heron (1972) for
matroids representable over fields, and generalized by Rota and Welsh in
the 1970s.

The conjecture generalizes Read's conjecture (since the chromatic
polynomial of a graph is a special case of the characteristic polynomial
of a matroid). It also includes, as a special case, the unimodality of
the coefficients of the \textbf{Tutte polynomial} and related
combinatorial invariants.

\subsection{4.2 Progress in Stages}\label{ch2022junehuh:progress-in-stages}

Huh resolved the conjecture in stages, each stage requiring deeper and
more powerful geometric techniques:

\textbf{Stage 1 (2012):} Huh proved the conjecture for matroids
representable over fields of characteristic zero, extending his work on
graphs. This used the theory of \textbf{hyperplane arrangements} and the
cohomology of \textbf{reciprocal planes}.

\textbf{Stage 2 (2012, with Eric Katz):} Huh and Katz extended the proof
to matroids representable over fields of arbitrary characteristic, using
\textbf{tropical geometry} and the theory of \textbf{Bergman fans}. The
Bergman fan of a matroid is a polyhedral fan that encodes the matroid's
structure; Huh and Katz showed that this fan can be interpreted as the
tropicalization of the complement of a hyperplane arrangement, allowing
the use of tropical intersection theory.

\textbf{Stage 3 (2015, with Karim Adiprasito and Eric Katz):} The full
conjecture was proved for \textbf{all matroids}, including those not
representable over any field. This required an entirely new approach:
the construction of a \textbf{Hodge theory for matroids} that does not
rely on any underlying geometry but is intrinsic to the matroid itself.
The proof was published in the \emph{Annals of Mathematics} in 2018.

\subsection{4.3 Impact}\label{ch2022junehuh:impact-1}

The resolution of the Heron--Rota--Welsh conjecture was recognized as
a significant achievement in combinatorics in decades. It
completed a program that had been pursued by some of the leading
combinatorialists of the 20th century. More importantly, the methods
developed for the proof --- particularly the Hodge theory of matroids
--- opened up entirely new directions in combinatorial research.

\begin{center}\rule{0.5\linewidth}{0.5pt}\end{center}

\section{5. Hodge Theory for Matroids}\label{ch2022junehuh:hodge-theory-for-matroids}

\subsection{5.1 What is Hodge Theory?}\label{ch2022junehuh:what-is-hodge-theory}

Classical Hodge theory, developed by W. V. D. Hodge in the 1930s and
refined by many others, is a set of deep results about the cohomology of
smooth projective algebraic varieties. The three pillars of Hodge theory
are:

\begin{enumerate}
\def\labelenumi{\arabic{enumi}.}
\tightlist
\item
 \textbf{The Hodge decomposition:} The cohomology of a smooth
 projective variety decomposes into \((p,q)\)-components.
\item
 \textbf{The hard Lefschetz theorem:} Multiplication by the class of a
 hyperplane section induces isomorphisms between certain cohomology
 groups.
\item
 \textbf{The Hodge--Riemann bilinear relations:} A certain bilinear
 form on cohomology is definite on the primitive components.
\end{enumerate}

These results have significant consequences, including the
\textbf{log-concavity of sequences} and the
\textbf{Brunn--Minkowski inequality}, and they imply
strong combinatorial constraints.

\subsection{5.2 Building Hodge Theory for
Matroids}\label{ch2022junehuh:building-hodge-theory-for-matroids}

The central achievement of Adiprasito--Huh--Katz (2015, published 2018)
was the construction of a commutative ring \(A^*(M)\) --- the
\textbf{Chow ring} of a matroid \(M\) --- together with an intersection
product satisfying analogues of the hard Lefschetz theorem and the
Hodge--Riemann relations.

The construction proceeds as follows:

\begin{itemize}
\tightlist
\item
 Given a matroid \(M\) on a ground set \(E\), one defines the
 \textbf{Bergman fan} \(\Sigma_M\), a polyhedral fan in
 \(\mathbb{R}^{|E|}\) that encodes the flats of the matroid.
\item
 The Chow ring \(A^*(M)\) is defined as the ring of piecewise
 polynomial functions on \(\Sigma_M\) modulo a certain ``linear''
 equivalence relation, analogous to the definition of the Chow ring of
 a toric variety.
\item
 A distinguished element \(L \in A^1(M)\) (the ``ample'' class) is
 chosen, corresponding to a choice of generic linear function.
\item
 Using techniques from \textbf{combinatorial Hodge theory} and the
 theory of \textbf{Minkowski weights}, the authors proved that
 multiplication by \(L\) satisfies the hard Lefschetz theorem and that
 the Hodge--Riemann bilinear relations hold.
\end{itemize}

This was completely unexpected. There is no underlying smooth projective
variety, no cohomology theory in the usual sense, and yet the abstract
combinatorial structure of the matroid is rich enough to support these
deep geometric properties.

\subsection{5.3 Consequences}\label{ch2022junehuh:consequences}

The Hodge--Riemann relations for matroids directly imply:

\begin{itemize}
\tightlist
\item
 \textbf{Log-concavity of the coefficients of the characteristic
 polynomial} (the Heron--Rota--Welsh conjecture).
\item
 \textbf{Log-concavity of the \(f\)-vector of the independence complex}
 (the Mason conjecture, in a stronger form).
\item
 \textbf{The Brunn--Minkowski inequality for matroids}, an analogue of
 the classical inequality from convex geometry.
\end{itemize}

\subsection{5.4 Impact}\label{ch2022junehuh:impact-2}

The development of Hodge theory for matroids has been described as ``one
of the most exciting developments in combinatorics in the 21st
century.'' It has opened the door to a new field ---
\textbf{combinatorial Hodge theory} --- in which Hodge-theoretic methods
are applied to purely combinatorial objects. Subsequent work by many
researchers has extended these ideas to polytopes, hyperplane
arrangements, and other combinatorial structures.

\begin{center}\rule{0.5\linewidth}{0.5pt}\end{center}

\section{6. The Dowling--Wilson Top-Heavy
Conjecture}\label{ch2022junehuh:the-dowlingwilson-top-heavy-conjecture}

\subsection{6.1 Background}\label{ch2022junehuh:background-2}

Consider a configuration \(P\) of \(n\) points spanning a
\(d\)-dimensional vector space. Let \(w_i\) be the number of linear
subspaces of dimension \(i\) spanned by the points. For example, if the
points are in general position, then \(w_i = \binom{n}{i+1}\) for
\(i < d\). The \textbf{Dowling--Wilson conjecture} (1974) asserts that
for any such point configuration, the sequence
\((w_0, w_1, \ldots, w_d)\) is \textbf{unimodal}:
\(w_0 \le w_1 \le \cdots \le w_{k} \ge \cdots \ge w_d\) for some \(k\).
In fact, the conjecture states something stronger: the sequence is
\textbf{top-heavy}, meaning that \(w_i \le w_{d-i}\) for \(i \le d/2\).

For point configurations in general position, the equality case is
known. But for arbitrary configurations, the conjecture had resisted
proof for nearly 50 years.

\subsection{6.2 Huh and Wang's Proof}\label{ch2022junehuh:huh-and-wangs-proof}

In 2017, Huh and Botong Wang proved the Dowling--Wilson conjecture using
methods from algebraic geometry. The key insight was to interpret the
numbers \(w_i\) as the dimensions of the cohomology groups of a certain
\textbf{wonderful compactification} of the complement of a hyperplane
arrangement.

A wonderful compactification, introduced by De Concini and Procesi, is a
smooth projective variety that compactifies the complement of a
hyperplane arrangement while preserving its combinatorial structure. The
cohomology of this variety is known to have a \textbf{Lefschetz
decomposition} with respect to a hyperplane class, which forces the
dimensions to satisfy the hard Lefschetz theorem.

By carefully analyzing the cohomology of the wonderful compactification,
Huh and Wang proved that the dimensions \(w_i\) satisfy the desired
inequalities. The proof used:

\begin{itemize}
\tightlist
\item
 The theory of \textbf{hyperplane arrangements} and their combinatorial
 structure.
\item
 The \textbf{cohomology of wonderful compactifications}, including the
 description of their cohomology rings.
\item
 The \textbf{hard Lefschetz theorem} for smooth projective varieties.
\end{itemize}

\subsection{6.3 Impact}\label{ch2022junehuh:impact-3}

The proof of the Dowling--Wilson conjecture was a major achievement that
further demonstrated the power of algebraic-geometric methods in
combinatorics. It provided a complete solution to a problem that had
been open for decades and had been studied by some of the leading
figures in combinatorics.

\begin{center}\rule{0.5\linewidth}{0.5pt}\end{center}

\section{7. Lorentzian Polynomials}\label{ch2022junehuh:lorentzian-polynomials}

\subsection{7.1 Background}\label{ch2022junehuh:background-3}

In 2020, Huh and Petter Branden introduced the theory of
\textbf{Lorentzian polynomials}, a new class of multivariate polynomials
with powerful combinatorial properties. A polynomial
\(f(x_1, \ldots, x_n)\) with nonnegative coefficients is called
\textbf{Lorentzian} if its Hessian matrix has at most one positive
eigenvalue at every point in the positive orthant. This definition is
inspired by the theory of Lorentzian metrics in differential geometry.

The class of Lorentzian polynomials includes many important families:

\begin{itemize}
\tightlist
\item
 \textbf{Stable polynomials} (multivariate polynomials with no roots in
 the positive orthant).
\item
 \textbf{Volume polynomials} of convex bodies.
\item
 \textbf{Elementary symmetric polynomials}.
\item
 \textbf{The basis generating functions of matroids}.
\end{itemize}

\subsection{7.2 Key Properties}\label{ch2022junehuh:key-properties}

Huh and Branden proved that Lorentzian polynomials satisfy strong
combinatorial properties:

\begin{enumerate}
\def\labelenumi{\arabic{enumi}.}
\tightlist
\item
 \textbf{Log-concavity:} The coefficients of a Lorentzian polynomial
 (when expanded in the appropriate basis) form a log-concave sequence.
\item
 \textbf{Closure properties:} The class of Lorentzian polynomials is
 closed under multiplication, convolution, and taking partial
 derivatives.
\item
 \textbf{Rayleigh monotonicity:} A certain technical condition that
 implies the \textbf{strong Rayleigh property}.
\end{enumerate}

\subsection{7.3 Applications}\label{ch2022junehuh:applications}

The theory of Lorentzian polynomials provides a unified framework for
proving log-concavity across a wide range of combinatorial structures:

\begin{itemize}
\tightlist
\item
 \textbf{Matroid basis generating functions:} The basis generating
 function of any matroid is Lorentzian, implying the log-concavity of
 the \textbf{\(f\)-vector} and the \textbf{\(h\)-vector} of the
 matroid.
\item
 \textbf{Graphical and Tutte polynomials:} Certain specializations are
 Lorentzian.
\item
 \textbf{Khovanskii--Pukhlikov inequalities:} The theory provides a new
 proof of these important inequalities from convex geometry.
\item
 \textbf{Symmetric function inequalities:} Many classical inequalities
 for Schur functions and other symmetric functions can be derived from
 Lorentzianity.
\end{itemize}

\subsection{7.4 Impact}\label{ch2022junehuh:impact-4}

The development of Lorentzian polynomials has been described as a
``significant advance'' for combinatorial inequalities. It provides a unified
explanation for hundreds of known log-concavity results and has become
an active area of research. The paper was published in the \emph{Annals
of Mathematics} in 2020 and has already received extensive citations.

\begin{center}\rule{0.5\linewidth}{0.5pt}\end{center}

\section{8. The Strong Mason
Conjecture}\label{ch2022junehuh:the-strong-mason-conjecture}

\subsection{8.1 Background}\label{ch2022junehuh:background-4}

The \textbf{Mason conjecture} (1972) concerns the \textbf{\(f\)-vector}
of the independence complex of a matroid. Let \(f_k\) be the number of
independent sets of size \(k\) in a matroid \(M\). The conjecture states
that the sequence \((f_0, f_1, \ldots, f_r)\) is \textbf{log-concave}:
\(f_k^2 \ge f_{k-1} f_{k+1}\) for all \(k\). A weaker version (the
``ultra-log-concavity'' conjecture) was also proposed by Mason.

These conjectures generalize the classical theorem that the binomial
coefficients are log-concave (the matroid of a projective space has
\(f_k = \binom{n}{k}\)).

\subsection{8.2 Huh's Proof}\label{ch2022junehuh:huhs-proof-1}

In 2015, Huh (as part of the Adiprasito--Huh--Katz work) proved the
strong Mason conjecture. The proof followed from the Hodge--Riemann
relations for matroids, using the interpretation of the Chow ring and
the hard Lefschetz theorem.

Specifically, the numbers \(f_k\) can be expressed as the coefficients
of the \textbf{Tutte polynomial} of the matroid evaluated at certain
parameters. Using the Hodge--Riemann relations for the matroid's Chow
ring, Huh and his coauthors proved that these coefficients satisfy the
strong Mason inequalities.

\subsection{8.3 Impact}\label{ch2022junehuh:impact-5}

The proof of the strong Mason conjecture completed a program that had
been pursued by matroid theorists for over 40 years. It demonstrated the
power of the Hodge-theoretic approach and provided further evidence that
the Chow ring of a matroid captures essential combinatorial information.

\begin{center}\rule{0.5\linewidth}{0.5pt}\end{center}

\section{9. Lagrangian Geometry of
Matroids}\label{ch2022junehuh:lagrangian-geometry-of-matroids}

\subsection{9.1 Overview}\label{ch2022junehuh:overview}

In a series of papers with Federico Ardila and Graham Denham
(2021--2023), Huh developed a \textbf{Lagrangian} approach to matroid
theory. The central object is a \textbf{Lagrangian matroid} --- a
matroid equipped with additional data analogous to a Lagrangian subspace
in symplectic geometry.

\subsection{9.2 Key Results}\label{ch2022junehuh:key-results}

The Lagrangian geometry of matroids provides:

\begin{itemize}
\tightlist
\item
 A new interpretation of the \textbf{basis exchange axiom} of matroid
 theory in terms of Lagrangian intersections.
\item
 A generalization of the \textbf{Hodge--Riemann relations} to the
 Lagrangian setting.
\item
 New inequalities for matroid invariants, including the
 \textbf{convexity of the Tutte polynomial}.
\item
 Connections to the theory of \textbf{Coxeter groups} and
 \textbf{Lorentzian polynomials}.
\end{itemize}

\subsection{9.3 Impact}\label{ch2022junehuh:impact-6}

This work extends the scope of combinatorial Hodge theory and
establishes new connections between matroid theory, symplectic geometry,
and combinatorial optimization. It provides a framework for proving
inequalities for matroids that were previously inaccessible.

\begin{center}\rule{0.5\linewidth}{0.5pt}\end{center}

\section{10. Stellahedral Geometry of
Matroids}\label{ch2022junehuh:stellahedral-geometry-of-matroids}

\subsection{10.1 Overview}\label{ch2022junehuh:overview-1}

In a 2023 paper with Christopher Eur and Matt Larson, Huh introduced the
\textbf{stellahedron} of a matroid --- a convex polytope associated to a
matroid that encodes its combinatorial structure. The stellahedron is
the \textbf{Minkowski sum} of simplices corresponding to the flats of
the matroid.

\subsection{10.2 Key Results}\label{ch2022junehuh:key-results-1}

The stellahedral geometry provides:

\begin{itemize}
\tightlist
\item
 A geometric model for the Chow ring of the matroid: the Chow ring is
 isomorphic to the ring of piecewise polynomial functions on the
 stellahedron.
\item
 A new proof of the hard Lefschetz theorem for matroids using the
 geometry of the stellahedron.
\item
 The \textbf{stellahedral volume} of a matroid, a new invariant that
 generalizes the volume of a convex polytope and satisfies the
 Hodge--Riemann relations.
\end{itemize}

\subsection{10.3 Impact}\label{ch2022junehuh:impact-7}

The stellahedral construction gives a concrete, visualizable geometric
model for the abstract Hodge theory of matroids. It has made the theory
more accessible and has led to new results on the combinatorial
applications of Hodge theory.

\begin{center}\rule{0.5\linewidth}{0.5pt}\end{center}

\section{11. The Semi-Small Decomposition and Intersection Cohomology of
Matroids}\label{ch2022junehuh:the-semi-small-decomposition-and-intersection-cohomology-of-matroids}

\subsection{11.1 Overview}\label{ch2022junehuh:overview-2}

In a 2022 paper with Tom Braden, Jacob Matherne, Nick Proudfoot, and
Botong Wang, Huh developed the \textbf{intersection cohomology} of
matroids. Intersection cohomology, introduced by Goresky and MacPherson,
is a refinement of ordinary cohomology that satisfies Poincare duality
even for singular spaces.

\subsection{11.2 Key Results}\label{ch2022junehuh:key-results-2}

The authors proved that the intersection cohomology of a matroid
(defined via the Bergman fan) satisfies:

\begin{itemize}
\tightlist
\item
 \textbf{Poincare duality:} A non-degenerate pairing between \(IH^i\)
 and \(IH^{2d-i}\) where \(d\) is the rank of the matroid.
\item
 \textbf{The hard Lefschetz theorem:} Multiplication by a distinguished
 class induces an isomorphism \(IH^{d-i} \to IH^{d+i}\).
\item
 \textbf{The Hodge--Riemann relations:} A positivity condition for the
 intersection pairing.
\item
 \textbf{A semi-small decomposition:} The intersection cohomology of a
 matroid decomposes into a direct sum of simple pieces, analogous to
 the decomposition theorem in algebraic geometry.
\end{itemize}

\subsection{11.3 Impact}\label{ch2022junehuh:impact-8}

This work completes the Hodge-theoretic picture for matroids by
establishing the full analogue of the decomposition theorem, one of the
fundamental results in algebraic geometry. The semi-small decomposition
provides an effective tool for computing matroid invariants and for
proving new inequalities.

\begin{center}\rule{0.5\linewidth}{0.5pt}\end{center}

\section{12. Combinatorial Applications of the Hodge--Riemann
Relations}\label{ch2022junehuh:combinatorial-applications-of-the-hodgeriemann-relations}

The Hodge--Riemann relations for matroids --- and the broader theory of
Lorentzian polynomials --- have found applications well beyond the
original conjectures that motivated them:

\subsection{12.1 Inequalities for Graph
Polynomials}\label{ch2022junehuh:inequalities-for-graph-polynomials}

The chromatic polynomial, the Tutte polynomial, and the reliability
polynomial of a graph all satisfy log-concavity and other inequalities
that follow from the Hodge--Riemann relations.

\subsection{12.2 Convex Geometry}\label{ch2022junehuh:convex-geometry}

The \textbf{volume polynomial} of a convex body (or, more generally, a
mixed volume) is Lorentzian. This gives a new proof of the classical
Aleksandrov--Fenchel inequalities for mixed volumes and provides a
bridge between matroid theory and convex geometry.

\subsection{12.3 Algebraic Statistics}\label{ch2022junehuh:algebraic-statistics}

Lorentzian polynomials have been applied to the study of \textbf{maximum
likelihood estimation} for discrete statistical models, where they imply
that the likelihood function has a unique maximum.

\subsection{12.4 Machine Learning and
Optimization}\label{ch2022junehuh:machine-learning-and-optimization}

The log-concavity properties implied by Lorentzianity are important in
\textbf{Bayesian optimization} and \textbf{convex optimization}, where
they guarantee that certain functions are well-behaved for numerical
algorithms.

\begin{center}\rule{0.5\linewidth}{0.5pt}\end{center}

\section{13. Broader Impact on
Mathematics}\label{ch2022junehuh:broader-impact-on-mathematics}

\subsection{13.1 Transformation of
Combinatorics}\label{ch2022junehuh:transformation-of-combinatorics}

Huh's work has changed the relationship between
combinatorics and algebraic geometry. What was once a one-way street ---
combinatorics providing examples for geometers --- has become a two-way
exchange, with deep geometric theorems being used to solve combinatorial
problems.

The field of \textbf{combinatorial Hodge theory}, essentially created by
Huh and his collaborators, is now an active area of research with dozens
of practitioners. Researchers are applying Hodge-theoretic methods to:

\begin{itemize}
\tightlist
\item
 \textbf{Polytopes and fans} (combinatorial analogues of toric
 varieties)
\item
 \textbf{Hyperplane arrangements} (connections to the topology of
 complements)
\item
 \textbf{Coxeter groups} (connections to the Bruhat order)
\item
 \textbf{Cluster algebras} (connections to the theory of positive
 matrices)
\item
 \textbf{Representation theory} (connections to the cohomology of flag
 varieties)
\end{itemize}

\subsection{13.2 Unification of
Mathematics}\label{ch2022junehuh:unification-of-mathematics}

Huh's work is an illustration of the unity of mathematics. By
showing that the same Hodge-theoretic structures that govern the
cohomology of algebraic varieties also govern the combinatorics of
matroids, he has revealed a deep structural parallelism between the
continuous world of geometry and the discrete world of combinatorics.

As Huh himself has said: ``There is an underlying unity between all the
different ways of doing mathematics, be it algebraic, analytical,
geometric, combinatorial, or another. My goal is to uncover this
underlying unity as much as possible and build a framework that allows
problems to be tackled with diverse methods.''

\subsection{13.3 Influence on Algebraic
Geometry}\label{ch2022junehuh:influence-on-algebraic-geometry}

Huh's work has also influenced algebraic geometry itself. The proof of
the Heron--Rota--Welsh conjecture required the development of new
techniques that have applications to the study of \textbf{tropical
varieties}, \textbf{wonderful compactifications}, and \textbf{the
geometry of hyperplane arrangements}. The construction of Chow rings for
matroids has led to a better understanding of the \textbf{cohomology of
toric varieties} and the \textbf{intersection theory of fans}.

\subsection{13.4 A New Generation of
Mathematicians}\label{ch2022junehuh:a-new-generation-of-mathematicians}

Huh's unconventional career path has made him a role model for
mathematicians who may not fit the traditional ``prodigy'' mold. His
story demonstrates that:

\begin{itemize}
\tightlist
\item
 Mathematical creativity can emerge at any stage of life.
\item
 A background in other fields (Huh studied physics and astronomy and
 wrote poetry) can provide unique perspectives.
\item
 Mentorship can be substantial (Hironaka's influence was decisive).
\item
 Failure and discouragement in early education do not preclude later
 success.
\end{itemize}

\subsection{13.5 Public Engagement}\label{ch2022junehuh:public-engagement}

Since winning the Fields Medal, Huh has become an important voice in
public discussions about mathematics. He has spoken about the importance
of creativity in mathematics, the value of interdisciplinary thinking,
and the need to support researchers who take unconventional paths. His
willingness to share his personal story has motivated many students.

\begin{center}\rule{0.5\linewidth}{0.5pt}\end{center}

\section{14. Unusual Career Trajectory: Lessons for the
Profession}\label{ch2022junehuh:unusual-career-trajectory-lessons-for-the-profession}

\subsection{14.1 A Fields Medalist Who Struggled with
Math}\label{ch2022junehuh:a-fields-medalist-who-struggled-with-math}

Huh's path to the Fields Medal is notable in many ways. He is
perhaps the first Fields Medalist to have dropped out of high school, to
have failed mathematics courses, and to have seriously pursued poetry as
an alternative career. His journey challenges common assumptions
about mathematical talent and its development.

\subsection{14.2 The Role of Mentorship}\label{ch2022junehuh:the-role-of-mentorship}

Huh's encounter with Heisuke Hironaka was substantial. Hironaka did
not simply teach Huh mathematics; he showed him what mathematics
\emph{is} --- a creative, human activity driven by curiosity and
discovery. This mentorship provided the emotional and intellectual
foundation for Huh's subsequent development.

\subsection{14.3 The Late Bloomer in
Mathematics}\label{ch2022junehuh:the-late-bloomer-in-mathematics}

Huh's story is a notable counterexample to the ``child prodigy''
narrative that dominates popular perceptions of mathematical talent.
While early exposure and aptitude can certainly be helpful, they are not
prerequisites for significant contributions. Huh's case suggests that
mathematical creativity can flourish when individuals find the right
environment, mentorship, and problems that match their interests.

\subsection{14.4 Implications for
Education}\label{ch2022junehuh:implications-for-education}

Huh's experience has implications for mathematics education:

\begin{itemize}
\tightlist
\item
 \textbf{Multiple pathways:} Routes to mathematical success are
 diverse; rigid tracking based on early performance may miss late
 bloomers.
\item
 \textbf{The importance of intellectual excitement:} Huh was drawn to
 mathematics not by grades or prizes but by the excitement of seeing
 research in progress.
\item
 \textbf{The value of broad preparation:} Huh's background in physics
 and his interest in poetry may have contributed to his unique
 perspective on mathematical problems.
\end{itemize}

\begin{center}\rule{0.5\linewidth}{0.5pt}\end{center}

\section{15. Awards and Recognition}\label{ch2022junehuh:awards-and-recognition}

\textbf{2022:} - Fields Medal (International Mathematical Union) -
MacArthur Fellowship (MacArthur Foundation)

\textbf{2021:} - Ho-Am Prize in Science (Samsung Foundation)

\textbf{2019:} - New Horizons in Mathematics Prize (Breakthrough Prize
Foundation)

\textbf{2018:} - Invited Speaker, International Congress of
Mathematicians (Rio de Janeiro)

\textbf{2017:} - Blavatnik Award for Young Scientists (US Regional)

\textbf{2016:} - Clay Research Fellow, Clay Mathematics Institute

\textbf{Earlier:} - National Academy of Sciences of Korea (Young
Scientist Award)

\begin{center}\rule{0.5\linewidth}{0.5pt}\end{center}

\section{16. Selected Major
Publications}\label{ch2022junehuh:selected-major-publications}

\begin{enumerate}
\def\labelenumi{\arabic{enumi}.}
\item
 \textbf{``Milnor numbers of projective hypersurfaces and the chromatic
 polynomial of graphs.''} \emph{Journal of the American Mathematical
 Society} 25 (2012), 907--927. --- \emph{Proof of Read's conjecture on
 unimodality of chromatic polynomial coefficients.}
\item
 \textbf{``The characteristic polynomial of a matroid''} (with E.
 Katz). \emph{Annals of Mathematics} 176 (2012), 2099--2136. ---
 \emph{Proof of Heron--Rota--Welsh conjecture for matroids
 representable over fields.}
\item
 \textbf{``Hodge theory of matroids''} (with K. Adiprasito and E.
 Katz). \emph{Annals of Mathematics} 188 (2018), 381--452. ---
 \emph{Proof of the full Heron--Rota--Welsh conjecture and the strong
 Mason conjecture via Hodge theory for matroids.}
\item
 \textbf{``Combinatorial applications of the Hodge--Riemann
 relations''} (with K. Adiprasito, E. Katz, and others).
 \emph{Proceedings of the ICM} (2018). --- \emph{Survey of applications
 of Hodge-theoretic methods to combinatorics.}
\item
 \textbf{``Lorentzian polynomials''} (with P. Branden). \emph{Annals of
 Mathematics} 192 (2020), 821--891. --- \emph{Development of the theory
 of Lorentzian polynomials.}
\item
 \textbf{``Logarithmic concavity for morphisms of matroids''} (with C.
 Eur). \emph{Advances in Mathematics} 367 (2020), 107094. ---
 \emph{Extension of Hodge theory to morphisms of matroids.}
\item
 \textbf{``Lagrangian geometry of matroids''} (with F. Ardila and G.
 Denham). \emph{Journal of the American Mathematical Society} 36
 (2023), 727--794. --- \emph{Development of Lagrangian matroid theory.}
\item
 \textbf{``Stellahedral geometry of matroids''} (with C. Eur and M.
 Larson). \emph{Forum of Mathematics, Pi} 11 (2023), 1--48. ---
 \emph{Geometric model for the Chow ring of a matroid.}
\item
 \textbf{``A semi-small decomposition of the Chow ring of a matroid''}
 (with T. Braden, J. Matherne, N. Proudfoot, and B. Wang).
 \emph{Advances in Mathematics} 409 (2022), 108646. ---
 \emph{Intersection cohomology and decomposition theorem for matroids.}
\item
 \textbf{``Combinatorics and Hodge theory.''} \emph{Proceedings of the
 International Congress of Mathematicians} 1 (2022), 212--239. ---
 \emph{Huh's own survey of his work.}
\end{enumerate}

\begin{center}\rule{0.5\linewidth}{0.5pt}\end{center}

\section{17. Ongoing Work and Future
Directions}\label{ch2022junehuh:ongoing-work-and-future-directions}

Huh continues to be highly active in research, with several ongoing
projects:

\subsection{17.1 Extending the Hodge-Theoretic
Framework}\label{ch2022junehuh:extending-the-hodge-theoretic-framework}

Working with collaborators, Huh is extending the Hodge theory of
matroids to other combinatorial structures, including \textbf{oriented
matroids}, \textbf{polymatroids}, and \textbf{convex geometries}. The
goal is to develop a general framework in which Hodge-theoretic
inequalities can be proved for any structure that satisfies certain
combinatorial axioms.

\subsection{17.2 Connections to Representation
Theory}\label{ch2022junehuh:connections-to-representation-theory}

Recent work has explored connections between the Hodge theory of
matroids and the representation theory of \textbf{Coxeter groups} and
\textbf{Hecke algebras}. These connections suggest that the
Hodge--Riemann relations for matroids may have deeper origins in the
algebraic combinatorics of reflection groups.

\subsection{17.3 Algorithmic
Applications}\label{ch2022junehuh:algorithmic-applications}

The theory of Lorentzian polynomials has potential applications in
\textbf{optimization} and \textbf{algorithmic combinatorics}. Huh has
expressed interest in exploring whether the log-concavity properties
implied by Lorentzianity can be exploited for efficient computation.

\subsection{17.4 Mentoring the Next
Generation}\label{ch2022junehuh:mentoring-the-next-generation}

Huh is actively involved in mentoring graduate students and postdocs at
Princeton and through the Korea Institute for Advanced Study. He is
committed to building a community of researchers working at the
interface of combinatorics and algebraic geometry.

\begin{center}\rule{0.5\linewidth}{0.5pt}\end{center}

\section{18. Conclusion}\label{ch2022junehuh:conclusion}

June Huh has reshaped the landscape of combinatorics by importing deep
ideas from algebraic geometry --- particularly Hodge theory --- into the
study of discrete structures. His proofs of the Heron--Rota--Welsh
conjecture, the Dowling--Wilson conjecture, and the strong Mason
conjecture resolved problems that had stood open for decades, and his
development of Lorentzian polynomials has provided the field with a
new tool whose applications continue to multiply.

What makes Huh's achievement notable is not only its mathematical
depth but its conceptual scope. By revealing that the same
Hodge-theoretic structures that govern the cohomology of algebraic
varieties also govern the combinatorics of matroids, he has uncovered a
deep unity in mathematics. His work has transformed combinatorics from a
field that was largely self-contained to one that draws on, and
contributes to, some of the major areas of algebraic geometry.

Huh's personal story --- from high school dropout and aspiring poet to
Fields Medalist --- has inspired mathematicians and non-mathematicians
alike. It serves as a powerful reminder that mathematical creativity is
not limited to those who show early promise, and that the path to the
highest levels of achievement can be as unconventional as the
mathematics itself.

The Fields Medal citation captures the breadth of his contributions:
bringing Hodge theory to combinatorics, proving the Dowling--Wilson
conjecture, proving the Heron--Rota--Welsh conjecture, developing
Lorentzian polynomials, and proving the strong Mason conjecture. Each of
these achievements alone would be a distinguished career; together, they
represent an influential body of work in 21st-century
mathematics.

\begin{center}\rule{0.5\linewidth}{0.5pt}\end{center}

\section{19. References and Further
Reading}\label{ch2022junehuh:references-and-further-reading}

\subsection{Official Sources}\label{ch2022junehuh:official-sources}

\begin{itemize}
\tightlist
\item
 IMU Fields Medal 2022 Announcement:
 https://www.mathunion.org/imu-awards/fields-medal/fields-medal-2022
\item
 Fields Medal Citation for June Huh:
 https://www.mathunion.org/fileadmin/IMU/Prizes/Fields/2022/IMU\_Fields22\_Huh\_citation.pdf
\item
 Kalai, G. ``The Work of June Huh.'' ICM 2022 Laudatio.
 https://www.mathunion.org/fileadmin/IMU/Prizes/Fields/2022/laudatio-jh.pdf
\end{itemize}

\subsection{Biographical and Overview
Articles}\label{ch2022junehuh:biographical-and-overview-articles}

\begin{itemize}
\tightlist
\item
 Hartnett, K. ``A Path Less Taken to the Peak of the Math World.''
 \emph{Quanta Magazine}, June 27, 2017.
\item
 Chang, K. ``Fields Medals in Mathematics Won by Four Under Age 40.''
 \emph{New York Times}, July 5, 2022.
\item
 Princeton University News: ``Princeton mathematician June Huh awarded
 prestigious Fields Medal.'' July 5, 2022.
\item
 Baker, M. ``A Fields Medal for June Huh.'' \emph{Matt Baker's Math
 Blog}, July 5, 2022.
\end{itemize}

\subsection{Technical Surveys}\label{ch2022junehuh:technical-surveys}

\begin{itemize}
\tightlist
\item
 Eur, C. ``Essence of independence: Hodge theory of matroids since June
 Huh.'' \emph{Bulletin of the AMS} (2023). arXiv:2211.05724.
\item
 Huh, J. ``Combinatorics and Hodge theory.'' \emph{Proceedings of the
 ICM} 1 (2022), 212--239.
\item
 Adiprasito, K., Huh, J., and Katz, E. ``Hodge theory of matroids.''
 \emph{Notices of the AMS} 64 (2017), 26--30.
\end{itemize}

\subsection{Background Reading}\label{ch2022junehuh:background-reading}

\begin{itemize}
\tightlist
\item
 Oxley, J. \emph{Matroid Theory}. 2nd ed., Oxford University Press,
 2011.
\item
 Griffiths, P. and Harris, J. \emph{Principles of Algebraic Geometry}.
 Wiley, 1978. (For classical Hodge theory.)
\item
 Maclagan, D. and Sturmfels, B. \emph{Introduction to Tropical
 Geometry}. AMS, 2015.
\item
 Bränden, P. and Huh, J. ``Lorentzian polynomials.'' \emph{Annals of
 Mathematics} 192 (2020), 821--891.
\end{itemize}
